\section{Phân tích và Thảo luận Đối chiếu}

\subsection{Mối quan hệ kỹ thuật giữa lý thuyết cơ sở và bài báo hiện đại}
Dù ra đời ở hai giai đoạn công nghệ hoàn toàn khác biệt, cả lý thuyết tổng hợp biểu cảm truyền thống trong giáo trình và khung phương pháp hiện đại của bài báo đều chia sẻ một nền tảng nhận thức chung: biểu cảm khuôn mặt là hệ quả trực tiếp của các chuyển động cơ bắp cục bộ. Nếu như trong các phương pháp vật lý kinh điển, sự tác động của cơ bắp được tính toán thông qua lực đàn hồi của mô hình lò xo và khối lượng, thì ở phương pháp hiện đại, chuyển động cơ được mã hóa thành FAUs độc lập. Hơn nữa, ý tưởng về việc điều khiển cục bộ các vùng trên khuôn mặt, vốn được đề cập sơ khai qua việc chia nhỏ lưới bề mặt ba chiều trong các phương pháp biến đổi hình thái, nay đã được số hóa và kiểm soát một cách khéo léo thông qua Decoupled Cross Attention trong latent space của mô hình Diffusion.

\subsection{Sự tiến hóa về mặt phương pháp học}
Sự khác biệt cốt lõi nhất nằm ở cách thức kiến tạo hình ảnh. Các kỹ thuật cũ như ánh xạ biểu cảm hình học hay sử dụng ảnh tỷ lệ biểu cảm phụ thuộc nặng nề vào các thuật toán biến đổi lưới và tính toán quang học đơn giản. Những kỹ thuật này thường vấp phải giới hạn về tính tự nhiên khi gặp phải những khác biệt lớn về cấu trúc khuôn mặt giữa nguồn và đích, dễ dẫn đến hiện tượng biến dạng vật lý hoặc rách lưới nội suy ở các vùng biên. Ngược lại, việc ứng dụng mô hình Diffusion trong phương pháp mới không tiến hành chỉnh sửa trực tiếp trên cấu trúc điểm ảnh hay lưới ba chiều, mà sinh ra toàn bộ hình ảnh mới thông qua quá trình khử nhiễu liên tiếp. Sự chuyển dịch từ thao tác hình học thuần túy sang không gian mô hình tạo sinh có điều kiện giúp hệ thống không chỉ giải quyết triệt để lỗi cấu trúc vật lý, mà còn tái hiện được những thay đổi quang học vô cùng phức tạp mà các phương pháp cũ không thể mô phỏng.